\section{Introduction}
\label{sec:intro}

The growth of audiovisual content has driven the development of methods capable of analysing, indexing and automatically retrieving visual information. Within this landscape, multimodal video--text retrieval has consolidated as a relevant area at the intersection of computer vision and natural language processing. An especially challenging scenario is that of \emph{egocentric video}, captured from a first-person perspective using head- or body-mounted cameras. Such recordings register natural interactions between hands, objects and the surrounding environment and introduce a set of \emph{egocentric challenges}---strong viewpoint, illumination and motion variability, dynamic camera motion, occlusions and fine-grained hand--object interactions---that models must accommodate in order to align visual content with textual descriptions.

The reference resource for this problem is the \href{https://epic-kitchens.github.io/}{EPIC-KITCHENS-100} dataset~\cite{damen2022rescaling}, which contains over 100\,h of egocentric kitchen video and approximately 90{,}000 narrated action segments. The associated \href{https://www.codabench.org/competitions/12008/\#/pages-tab}{Multi-Instance Retrieval (MIR) challenge} formalises the task: given a textual query, the system must rank 9{,}668 trimmed video segments from the test split (and vice versa) by their semantic relevance, where ground-truth relevance is provided as a continuous \emph{soft-label matrix} rather than as one-to-one binary correspondences~\cite{damen2022rescaling,wray2021semantic}. This formulation acknowledges that several segments may match a query at different degrees of similarity and so, is evaluated with order-sensitive metrics like \emph{normalized discounted cumulative gain} (nDCG)~\cite{jarvelin2002ndcg} and \emph{mean average precision} (mAP) rather than with the conventional Recall@K used in third-person retrieval. The egocentric challenges above, combined with the fact that most existing losses do not exploit the graded relevance signal provided by the challenge, make this benchmark hard for off-the-shelf multimodal models.

The EK-100 MIR experiments comprise three configurations evaluated under a unified Codabench protocol. The first is a \textbf{JPoSE-based zero-fine-tuning baseline}~\cite{wray2019jpose}, in which verb, noun and action sub-embeddings of pre-trained Part-of-Speech encoders are concatenated to produce a $9668\times 3842$ similarity matrix, reaching \textbf{53.53\,nDCG} on the official Codabench leaderboard. The second is a \textbf{score-level ensemble} that fuses three pre-trained retrieval models (JPoSE, MI-MM~\cite{miech2020milnce} and MMEN~\cite{wray2019jpose}) through row-wise min-max normalization followed by equal-weight averaging, optionally combined with a \textbf{graph-diffusion re-ranking} step that propagates similarity scores through visual and textual k-nearest-neighbour graphs in the spirit of~\cite{iscen2017diffusion}, reaching \textbf{55.82\,nDCG}. The third is an \textbf{end-to-end fine-tuning} of an AVION ViT-L/14 dual encoder~\cite{zhao2023avion} with the Symmetric Multi-Similarity (SMS) loss~\cite{wang2024symmetricmultisimilaritylossepickitchens100}, which consumes the soft-label relevancy matrices directly during training, complemented by a \emph{horizontal-flip + 32-frame} test-time augmentation strategy; this configuration, trained for sixteen epochs on a single GPU, reaches \textbf{69.68\,nDCG}---a \textbf{+16.15} absolute gain over the JPoSE baseline.
